% Chapter 3 -- Larger programs (model: TP 6.0 UG Ch 3 "Turbo Pascal
% units"; Turbo Initech has no units, ADR-0007 DEC-02, so the chapter
% teaches how a single-file program is organized).  AN-19.
% Sources (Law 1): os/tps/tps.pas ParseProgram (repeated, interleaved
% sections), the declare-before-use rule (os/tps/SYMBOLS-DUMP.md s7: a call
% to a later routine needs a prior forward); ADR-0007 DEC-02/03 idioms.
% Examples replayed by test-pubs-tps.
\Gates{test-tps-type-fpc, test-compiler, test-compiler-os, test-pubs-tps}
\chapter{Larger programs}
\ix{larger programs}
\ix{program structure}

A Turbo Initech program is always one file, however large it grows; there
are no units and no include files.  This chapter shows a complete program
of a few dozen lines and the habits that keep a larger program in order.

\section{A complete program}
The program \file{REPORTS.PAS} checks a batch of reports for cover sheets,
totals their pages, and records the count in a file:

\Program{reports}

\ProgramOutput{reports}

\section{The order of declarations}
\ix{declaration order}
Turbo Initech reads the program once, from top to bottom, and accepts a
name only after its declaration.  \file{REPORTS.PAS} is arranged so:

\begin{enumerate}
\item the constants \id{MaxReports} and \id{Missing}, which the declarations
  after them use;
\item the record type \id{Report};
\item the global variables that the routines share;
\item the routines, each after everything it uses;
\item a second variable section, for the variable \id{Bad} that only the
  main block uses;
\item the main block.
\end{enumerate}

Constant, type and variable sections may be repeated as often as you need,
in any order that keeps each declaration ahead of its uses.

\section{Routines that call each other}
\ix{forward declarations}
A routine may call only routines declared above it.  This program is
rejected, because \id{First} calls \id{Second} before \id{Second} is
declared:

\Source{err-later}

\CompilerOutput{err-later}{ERR}

The cure is a \rw{forward} declaration: the heading of \id{Second}, followed
by \rw{forward}, placed before \id{First}.  The full declaration comes later:

\Program{later}

\ProgramOutput{later}

\section{Techniques}
\SideHead{Searching}A search that may run off the end of an array stops on a
Boolean flag, and tests the index before it reads the element:

\Program{search}

\ProgramOutput{search}

\SideHead{Sorting}An insertion sort takes each element in turn and moves
the larger elements before it one place up until the place for it is
found.  The inner loop uses the same flag and the same nested test:

\Program{sort}

\ProgramOutput{sort}

\SideHead{Counting}An array indexed by \code{Ord(C) - Ord('A')} counts
letters; the two tests on \id{Code} keep the index within the array:

\Program{letters}

\ProgramOutput{letters}

\SideHead{Named states}Integer constants give names to the states of a
program, where another language would use an enumerated type:

\Program{states}

\ProgramOutput{states}

\section{A larger application: the staff directory}
\ix{staff directory}
\ix{application, sample}
The program \file{DIRECTRY.PAS} keeps a staff directory.  It saves the
staff to a file as lines of text, reads the file back into its tables,
sorts the people by name and prints a report.  It uses every technique of
this chapter, and it shows how a Turbo Initech program does without the
features the language leaves out.

\Program{directory}

On InitechDOS it writes \file{STAFF.TXT} and prints:

\ProgramOutput{directory}

\SideHead{Names in one string}A string cannot be an element of an array or a
field of a record.  The program therefore keeps all the names in the one
string \id{Pool}, \id{NameLen} characters each, and \id{GetName} and
\id{PutName} copy one name out of or into its slot.  The main block first
fills \id{Pool} with blanks so that every slot exists.

\SideHead{Records and an index}The badge and desk numbers are fields of the
records in \id{Staff}.  \id{Names} holds the slot numbers in the order the
report wants; \id{SortByName} rearranges this index, not the records, so
the records and the names stay in step.

\SideHead{Text in a file}\id{Save} writes each person as one line, the three
values separated by commas and the line ended by a line feed (\code{Chr(10)}).
\id{ReadLine} reads the file one byte at a time until a line feed, and
\id{Field} takes a number off the front of a line.  A text file written in
this way can be read with \cmd{TYPE} at the InitechDOS prompt.

\SideHead{One file, two uses}The file variable \id{Data} is global, because
a file cannot be passed as a parameter.  \std{assign}{Assign} names it once;
\id{Save} opens it with \std{rewrite}{Rewrite} and \id{Load} reopens it with
\std{reset}{Reset}, which closes it first.

\SideHead{Guards}Every loop that might run past the end of a string or an
array stops on a Boolean flag and tests the index in an \rw{if} of its own:
\id{ReadLine}, \id{Field} and the inner loop of \id{SortByName}.

\SideHead{Columns}\id{Report} pads each name with spaces to 13 characters,
since \std{write}{Write} has no field widths.

\section{Keeping a program in order}
\SideHead{Constants for sizes}Declare every array bound and string capacity
as a constant, as \id{MaxReports} is, so that one change resizes the
program.

\SideHead{Local variables}Declare a variable inside the routine that uses
it, as \id{SaveCount} declares \id{Log}.  Local variables exist only during
the call, and they cannot be changed by mistake elsewhere.

\SideHead{Guards}Turbo Initech evaluates both sides of \rw{and} and
\rw{or}.  When one test protects another, such as a bound check before an
array access, write two nested \rw{if} statements (the
\doctitle{Programmer's Guide}, Chapter 5).

\SideHead{Lists without pointers}Keep items in arrays and link them by
their indexes, with a constant such as \id{Missing} for ``none'' (the
\doctitle{Programmer's Guide}, Chapter 3).

\SideHead{Capacities}The compiler holds the whole symbol table of a program
at once.  Chapter 8 lists its capacities; a very large program may need
fewer or shorter names.

\section{Where to go from here}
Programs do not always do what you meant.  Chapter 4 shows you how to find
out why.
